\section{Configuration recipes}\label{sec:recipes}

The recipes below use only keys that exist in the shipped bundle (version \BundleVersion). The first listing is the
orchestrator configuration exactly as it ships; the others change one or two keys of it.

\paragraph{The shipped default (Fable, Opus and Sonnet hosts alike)} The same configuration serves every host: the
price gate keeps Opus sessions on Opus, \code{by\_host} sets Fable's own effort to medium, and \code{keep\_on\_host}
keeps review- and explain-like sessions on the host.\par
{\small
\lstinputlisting[basicstyle=\ttfamily\footnotesize, columns=fullflexible, frame=single, rulecolor=\color{okgray}]{generated/recipe-shipped.yaml}}

\paragraph{Recipe: accept the review and explain risk for the larger saving (Fable)} Remove the opt-out. Expected
cost falls from about \KpProxyTwo\X\ to about \KpRstarTwo\X\ of plain Fable, at the quality risk H7 measured.
\begin{lstlisting}[basicstyle=\ttfamily\footnotesize, frame=single, rulecolor=\color{okgray}]
model_routing:
  keep_on_host:
    task_types: []
\end{lstlisting}

\paragraph{Recipe: medium effort on a Sonnet host} Sonnet at medium cost \EcRatio\X\ its default. The per-host map
takes any model id:
\begin{lstlisting}[basicstyle=\ttfamily\footnotesize, frame=single, rulecolor=\color{okgray}]
effort_routing:
  by_host:
    claude-sonnet-5:
      strong: medium
\end{lstlisting}
Do not add Opus to this map: medium effort did not demonstrate a saving on Opus (\SHFourRatio\X).

\paragraph{Recipe: use the decision model instead of the rule} Jev chose the same session model as the rule on all
but \SHTwoDisc\ of \SHTwoReps\ scenario-reps, so this buys nothing measurable and sends the starting state to an
external service. If you want it anyway, opt in explicitly:
\begin{lstlisting}[basicstyle=\ttfamily\footnotesize, frame=single, rulecolor=\color{okgray}]
allow_external_state: true
model_routing:
  start_policy: judge
\end{lstlisting}

\paragraph{Claude Code, Codex and Copilot CLI} The smart tool makes the same session-start decision outside the
harness. \code{decide} prints the choice; \code{launch} makes it and starts the third-party agent with that model and
effort (\opt{dry-run} prints the command instead):
\begin{lstlisting}[basicstyle=\ttfamily\footnotesize, frame=single, rulecolor=\color{okgray}]
npx skills add michaeljabbour/amplifier-bundle-fast-decisions
amplifier-fast-decisions decide \
    --task "Fix the failing date parser test" \
    --host-model claude-fable-5-1 --workspace . --decider rules
amplifier-fast-decisions launch --harness claude \
    --task "Explain the cache layer" \
    --host-model claude-fable-5-1 --workspace . --dry-run
amplifier-fast-decisions launch --harness codex \
    --task "..." --host-model claude-fable-5-1
amplifier-fast-decisions launch --harness copilot \
    --task "..." --host-model claude-fable-5-1
npx skills update
\end{lstlisting}
Inside the Amplifier harness itself, add the bundle and keep it current with \code{amplifier bundle update}. Savings
outside the harness have not been measured yet.
